\documentclass[10pt,a4paper]{amsart}
\usepackage[margin=19mm]{geometry}
\usepackage{amsmath,amssymb,mathtools}
\usepackage[scr=rsfs,cal=euler]{mathalfa}
\usepackage{xfrac}
\usepackage[unicode,hidelinks,bookmarks=false]{hyperref}
\newcommand{\ie}{i.e.\ }
\newcommand{\cf}{cf.\ }
\newcommand{\resp}{respectively\ }
\newcommand{\Subsection}{\subsection}
\providecommand{\nobreakdash}{\nobreak\hbox{-}}
\setlength{\emergencystretch}{2em}
\multlinegap0pt
\usepackage{bm}
\usepackage{bbm}
\usepackage{dsfont}
\usepackage{xspace}
\usepackage{cancel}
\usepackage{mathtools}
\usepackage{amsthm}
\makeatletter
\@ifundefined{theorem}{\newtheorem{theorem}{Theorem}}{}
\makeatother
\newcommand{\F}{\mathbb{F}}
\newcommand{\Z}{\mathbb{Z}}
\title[Good Integers: (T,k)-Subclasses and Applications to Galois D]{Good Integers: (T,k)-Subclasses and Applications to Galois Duality in Coding Theory}
\author{Somphong Jitman, Panthakan Boonsuriyatham}
\date{Source: arXiv:2605.27933v1 (CC BY 4.0)}
\begin{document}
\maketitle

\noindent\emph{Source and licence.} Good Integers: (T,k)-Subclasses and Applications to Galois Duality in Coding Theory. Version arXiv:2605.27933v1 (CC BY 4.0), distributed under the Creative Commons Attribution 4.0 International licence, as recorded in the arXiv licence metadata for this exact version and shown on the abstract page https://arxiv.org/abs/2605.27933v1. Full rights notice: licenses/arxiv-2605.27933-CC-BY-4.0.txt.\par\smallskip

\noindent\textit{Editorial scope.} Counted unit: the Theorem whose source label is \texttt{thm:lcd-rr} in the article (source0.tex lines 2338--2349, source label \texttt{thm:lcd-rr}), delivered together with the article's own complete original proof of it (source0.tex lines 2351--2451, one \texttt{proof} environment of 101 lines). No mathematics is written, repaired or re-derived: statement and proof are the article's own text line for line. Required context retained uncounted: eq-LCD (lines 2045--2047); thm:reciprocal-coset (lines 2078--2091). Every definition, display and label of those lines is retained unchanged. Omitted from this package and not redistributed: 1--2044, 2048--2077, 2092--2337, 2350--2350, 2452--2894. Nothing inside a retained excerpt is omitted or altered: every retained body is delivered exactly as the article's lines are written, so no omission receipt of any kind accompanies this package. Citations: the retained excerpts carry no citation command at all, so no bibliography entry of the article is transcribed and no reference is redistributed. Reference adaptation: none was needed and none was made; every reference inside the delivered unit resolves inside it, so no unresolved reference or citation is produced. Printed slips kept literal: no printed word, symbol, hypothesis or number of the counted statement or of its proof was altered. Layout and class: the article's own class is \texttt{sn-jnl} with packages including graphicx, multirow, amsmath, amssymb, amsfonts, amsthm, mathrsfs, appendix, xcolor, textcomp, manyfoot, booktabs, algorithm, algorithmicx; the wrapper is the lane's pinned \texttt{amsart} editorial wrapper at 10pt on A4, which restates the theorem-like environments the retained text needs and adds the article's own macro definitions that the retained text needs (\texttt{\textbackslash F}, \texttt{\textbackslash Z}), with extra wrapper packages (bm, bbm, dsfont, xspace, cancel, mathtools). The delivered numbering is the unit's own. Assets: no retained span contains a figure environment, a graphics-inclusion command, a PDF-page inclusion, a table or a TikZ picture; no image file is copied into this package, the package declares no asset dependency and no image file is redistributed with it. Licence: version arXiv:2605.27933v1 of this article is distributed under the Creative Commons Attribution 4.0 International licence, as recorded in the arXiv licence metadata for this exact version and shown on the abstract page https://arxiv.org/abs/2605.27933v1. A full rights notice is delivered at licenses/arxiv-2605.27933-CC-BY-4.0.txt. Characters: every retained excerpt is pure ASCII, so the delivered engine, XeLaTeX with its default Latin Modern text font, reports no missing character.
\begin{align}
  \gcd(g(x),  h^{\ast,\theta}(x) )=1\label{eq-LCD},
\end{align}

\begin{theorem}\label{thm:reciprocal-coset}
For every $q^k$-cyclotomic class $Q\subseteq \Z_n$,
\[
f_Q^{\ast,\theta}(x)=f_{-q^TQ}(x),
\]
where
$
-q^TQ=\{-q^Ti \bmod n : i\in Q\}$.
Consequently, $f_Q$ is $\theta$-self-reciprocal if and only if
\[
Q=-q^TQ,
\]
or equivalently, if and only if every element of $Q$ has additive order in $G_{(T,k)}(q,1)$.
\end{theorem}

\begin{theorem}\label{thm:lcd-rr}
A cyclic code $C=\langle g(x)\rangle$ of length $n=n_0p^r$ over $\F_{q^k}$ is Galois LCD if and only if
\begin{align}\label{gen-LCD}
g(x)=
\left(\prod_{Q\in J_1} f_Q(x)^{p^r}\right)
\left(\prod_{Q\in J_2} F_Q(x)^{p^r}\right)
\end{align}
for some subsets $J_1\subseteq \Omega_1$ and $J_2\subseteq \Omega_2$. Consequently, the number of Galois LCD cyclic codes of length $n$ over $\F_{q^k}$ is
\[
2^{|\Omega_1|+|\Omega_2|}.
\]
\end{theorem}

\begin{proof}
Write
\[
g(x)=
\left(\prod_{Q\in \Omega_1} f_Q(x)^{a_Q}\right)
\left(\prod_{Q\in \Omega_2}\ \prod_{R\in \mathcal O(Q)} f_R(x)^{a_R}\right),
\]
where $0\le a_Q,a_R\le p^r$.
Then
\[
h(x)=\frac{x^n-1}{g(x)}
=
\left(\prod_{Q\in \Omega_1} f_Q(x)^{p^r-a_Q}\right)
\left(\prod_{Q\in \Omega_2}\ \prod_{R\in \mathcal O(Q)} f_R(x)^{p^r-a_R}\right).
\]
Since each $Q\in \Omega_1$ is fixed by $\sigma$,  by Theorem~\ref{thm:reciprocal-coset}
\[
h^{\ast,\theta}(x)=
\left(\prod_{Q\in \Omega_1} f_Q(x)^{p^r-a_Q}\right)
\left(\prod_{Q\in \Omega_2}\ \prod_{R\in \mathcal O(Q)} f_R(x)^{p^r-a_{\sigma^{-1}(R)}}\right).
\]
By \eqref{eq-LCD}, the code $C$ is $\theta$-LCD if and only if
$
\gcd\bigl(g(x),h^{\ast,\theta}(x)\bigr)=1$.
Since the irreducible factors $f_R(x)$ are pairwise coprime, this holds if and only if
$
\min\{a_R,\;p^r-a_{\sigma^{-1}(R)}\}=0$ for every $R$.
Equivalently,
\[
a_R=0
\quad\text{or}\quad
a_{\sigma^{-1}(R)}=p^r
\]
for every $R$.

We now consider the two parts separately.
First, let $Q\in \Omega_1$. Since $\sigma(Q)=Q$, the above condition becomes
$
a_Q=0 \text{ or }
a_Q=p^r$.
Hence, each factor $f_Q(x)$ with $Q\in \Omega_1$ occurs either with multiplicity $0$ or with multiplicity $p^r$.

Next, let $Q\in \Omega_2$, and write
$
\mathcal O(Q)=\{Q_0,Q_1,\dots,Q_{m-1}\}
$,  $ \sigma(Q_i)=Q_{i+1}$
 where the indices are computed modulo $m$.
Then the condition above becomes
$
a_{Q_i}=0 $ or $
a_{Q_{i-1}}=p^r$ 
 for all $i\in \{0,1,\dots, m-1\}$.
If all exponents on the orbit are zero, then the orbit contributes nothing to $g(x)$. Suppose now that at least one exponent is $
a_{Q_i}>0$
for some $i$. Then necessarily
$
a_{Q_{i-1}}=p^r>0$. Using the same argument, we conclude that $
a_{Q_{i-2}}=p^r$.
Repeating this process around the orbit, we obtain
\[
a_{Q_0}=a_{Q_1}=\cdots=a_{Q_{m-1}}=p^r.
\]
Therefore, for each orbit represented by $Q\in \Omega_2$, either all exponents are $0$ or all exponents are $p^r$. Hence, the contribution of such an orbit is either $1$ or
$
F_Q(x)^{p^r}$, where 
$F_Q(x)=\prod_{R\in \mathcal O(Q)} f_R(x)$.
It follows that
\[
g(x)=
\left(\prod_{Q\in J_1} f_Q(x)^{p^r}\right)
\left(\prod_{Q\in J_2} F_Q(x)^{p^r}\right)
\]
for some subsets $J_1\subseteq \Omega_1$ and $J_2\subseteq \Omega_2$.

Conversely, assume that
\[
g(x)=
\left(\prod_{Q\in J_1} f_Q(x)^{p^r}\right)
\left(\prod_{Q\in J_2} F_Q(x)^{p^r}\right)
\]
for some subsets $J_1\subseteq \Omega_1$ and $J_2\subseteq \Omega_2$. Then
\[
h(x)=
\left(\prod_{Q\in \Omega_1\setminus J_1} f_Q(x)^{p^r}\right)
\left(\prod_{Q\in \Omega_2\setminus J_2} F_Q(x)^{p^r}\right).
\]
Since each $f_Q(x)$ with $Q\in\Omega_1$ and each $F_Q(x)$ with $Q\in\Omega_2$ is $\theta$-self-reciprocal, we have
\[
h^{\ast,\theta}(x)=
\left(\prod_{Q\in \Omega_1\setminus J_1} f_Q(x)^{p^r}\right)
\left(\prod_{Q\in \Omega_2\setminus J_2} F_Q(x)^{p^r}\right).
\]
Thus, $g(x)$ and $h^{\ast,\theta}(x)$ have no common irreducible factor, which implies that $
\gcd\bigl(g(x),h^{\ast,\theta}(x)\bigr)=1$.
By \eqref{eq-LCD}, the code $C$ is $\theta$-LCD.

Finally, each factor indexed by $\Omega_1$ and each orbit product indexed by $\Omega_2$ can be chosen independently. Hence, the total number of $\theta$-LCD cyclic codes is
\[
2^{|\Omega_1|+|\Omega_2|}
\]as desired.
\end{proof}

\end{document}
